\chapter{TypeScript + p5.js を始める}
\label{chap:introduction}

この章では、TypeScript + p5.js のプログラムを自分のWindows環境で動かすための準備を行います。まだTypeScriptの詳しい文法やp5.jsの仕組みは深く説明しません。まずは、教材を取得し、プログラムを実行し、少し書き換え、Gitで保存するところまでを体験します。

\section{はじめに編}
ここでは、これから使うTypeScript、p5.js、Node.js、Gitがどのような役割を持つかを確認します。実行手順を暗記するのではなく、最初に全体像をつかむための編です。

\begin{chapterphasebox}{最初に読む}{blue}
この編は、教材を使い始めるときに一度読みます。用語が分からなくなったときは、あとから必要なところだけ見返してください。
\end{chapterphasebox}

\subsection{この章で行うこと}

プログラミングを学ぶときには、プログラムを書くためのエディタ、教材を取得するためのGit、プログラムを動かすためのNode.jsが必要になります。この章では、それらを順番に準備します。

\begin{table}[htbp]
    \centering
    \caption{第1章の3つの編と利用するタイミング}
    \label{tab:chapter01-three-phases}
    \begin{tabular}{p{0.17\linewidth}p{0.46\linewidth}p{0.25\linewidth}}
        \toprule
        編 & 主な内容 & 利用するタイミング \\
        \midrule
        はじめに編
            & TypeScript、p5.js、Node.js、Gitの役割
            & 最初に読み、必要なときに見返す \\
        準備編
            & インストール、CLI操作の練習、教材の取得、\texttt{npm install}
            & PC・教材ごとに原則1回 \\
        プログラムを動かす編
            & 開発サーバ、編集、実行、Gitのコミット
            & 授業や制作のたびに使う \\
        \bottomrule
    \end{tabular}
\end{table}

\begin{pointbox}{この章の目標}
この章の目標は、TypeScript + p5.js のプログラムをブラウザで表示できる状態にすることです。文法をすべて理解する必要はありません。まずは「動いた」という状態を作ることを優先します。
\end{pointbox}

\subsection{JavaScriptとTypeScriptの関係}

この教材で使う言語はTypeScriptです。
TypeScriptはJavaScriptをベースに、「型（データの種類の情報）」を追加した言語です。
TypeScriptは、開発中に書く言語で、実行するときにはJavaScriptに変換されます。
型情報を付けることで、間違った値を渡したときに実行前に気づきやすくなります。
JavaScriptはブラウザやNode.jsでそのまま動く言語で、Web開発では最も広く使われています。
JavaScriptはそのままでも十分強い言語ですが、規模の大きいプログラムでは次のような悩みが出やすくなります。

\begin{itemize}
    \item 関数にどんな値を渡すのかが分かりづらい
    \item 変数の値がいつどう変わるのか追いにくい
    \item チームで分担して開発すると、読み手ごとの解釈差が起きる
\end{itemize}

TypeScriptはこれらに対して、型注釈\footnote{英語ではtype annotationと呼びます。}
を通して「この値は何を表すのか」を明示しやすくします。
実行時にはTypeScriptを一度JavaScriptに変換しますが、書くときは型の手掛かりで理解しやすいコードになります。

近年TypeScriptが広く使われている理由は、JavaScriptの開発現場で「動きはそのままに、変更に強い構造」を求める流れが強まったためです。教材の中でも、変数・関数・クラスなどの概念を型付きで学ぶことで、次の章からのプログラミングがやりやすくなることを目指します。

\subsection{p5.jsとは何か}

p5.jsは、JavaScriptで図形、色、アニメーション、マウスやキーボードの操作などを扱いやすくするためのライブラリです。
この教材では、TypeScriptでプログラムを書き、p5.jsの機能を使って画面に図形を描いたり、動きのある作品を作ったりします。

p5.jsは、メディアアート系のプログラムを作るための言語であるProcessingの影響を受けて作られました。
Webブラウザで動くJavaScriptライブラリとして、図形を描き、動かし、操作に反応するプログラムを作るための機能を提供します。

p5.jsを利用して作ったプログラムのことを、よく\textbf{スケッチ}と呼びます。
スケッチという言葉には、絵やアイデアを試しながら形にしていくという意味があります。
この教材でも、p5.jsで作るひとまとまりのプログラムをスケッチと呼ぶことがあります。

\begin{pointbox}{p5.jsの背景}
p5.jsは、JavaScriptに図形の描画や入力の機能を加えるライブラリです。プログラムの中でp5.jsを読み込むと、これらの機能を使えます。
\end{pointbox}

\subsection{Node.jsとは何か}

Node.jsは、ブラウザの外でJavaScriptを動かすための実行環境です。
もともとJavaScriptはWebブラウザの中で動く言語として使われてきました。
しかし、TypeScriptをJavaScriptへ変換したり、開発用のサーバを起動したりする作業では、ブラウザの外でもJavaScriptの仕組みを使う必要があります。
そのために使うのがNode.jsです。

この教材では、Node.jsそのものを詳しく学ぶわけではありません。
Node.jsは、TypeScript + p5.jsの教材プロジェクトを動かすための土台として使います。
たとえば、\texttt{npm install}で必要なライブラリを準備したり、\texttt{npm run dev}でブラウザ確認用の開発サーバを起動したりします。

\begin{pointbox}{Node.jsとnpm}
Node.jsをインストールすると、\texttt{npm}というコマンドも一緒に使えるようになります。
\texttt{npm}は、p5.jsやTypeScriptなど、プロジェクトに必要なライブラリを準備するための道具です。
\end{pointbox}

\subsection{Gitとは何か}

Gitは、ファイルの変更履歴を記録するための道具です。
プログラムを書いていると、最初はうまく動いていたのに、少し書き換えたあとで動かなくなることがあります。
そのようなとき、Gitで途中の状態を記録しておくと、「どこで何を変更したのか」をあとから確認できます。

Gitは、単にファイルを保存する道具ではありません。
VS Codeでファイルを保存すると、現在の内容が上書きされます。
一方、Gitでは「この時点の変更を記録する」という形で、作業の節目を残します。
この記録を\textbf{コミット}と呼びます。

\begin{pointbox}{Gitでできること}
Gitを使うと、変更したファイルの一覧を確認したり、作業の節目を記録したりできます。
また、GitHubのようなサービスと組み合わせると、教材ファイルの配布や、複数のPCでの作業もしやすくなります。
\end{pointbox}

最初のうちは、Gitの仕組みをすべて理解する必要はありません。
この章では、次の3つの流れを覚えれば十分です。
\texttt{git status}で状態を確認し、\texttt{git add}で保存対象を選び、\texttt{git commit}で変更の記録を作ります。

Gitの考え方は少し抽象的なので、動画で全体像を見ておくのも役に立ちます。
次の動画は、Gitの導入として比較的見やすい内容です。
\url{https://www.youtube.com/watch?v=6SLMB7BPG9E}

もう少し詳しい説明として、次の動画も参考になります。
学生にとって少し難しく感じるかもしれませんが、Gitが何を解決する道具なのかを知る助けになります。
\url{https://www.youtube.com/watch?v=WMIiPcgGC4Q}

\section{準備編}
ここでは、教材を自分のPCで動かすための道具とファイルを準備します。多くの操作は、PCや教材ごとに原則として最初の1回だけ行います。

\begin{chapterphasebox}{PC・教材ごとに原則1回}{green}
すでに準備が終わっている場合、この編を授業のたびに繰り返す必要はありません。別のPCを使うときや、教材を取得し直すときに参照します。
\end{chapterphasebox}

\subsection{VS Codeをインストールする}

VS Codeは、プログラムを書くためのエディタです。TypeScriptのプログラムを読みやすく表示したり、間違いを見つけやすくしたりする機能があります。

\worktiming{PCごとに原則1回}

Webブラウザで\href{https://code.visualstudio.com/}{Visual Studio Codeの公式サイト}を開き、Windows用のインストーラをダウンロードします。図\ref{fig:chapter01-vscode-website}の「Download for Windows」が、そのためのボタンです。
この章の画面写真は操作の一例です。バージョン番号、表示言語、ボタンの位置などは、利用する時期やPCによって異なることがあります。

\begin{figure}[H]
    \centering
    \includegraphics[width=.95\linewidth]{figures/chapter01/vcode-oficial.png}
    \caption{VS Codeの公式サイトからWindows用のインストーラをダウンロードする}
    \label{fig:chapter01-vscode-website}
\end{figure}

ダウンロードが終わったら、Windowsのエクスプローラーで「ダウンロード」フォルダを開きます。
図\ref{fig:chapter01-vscode-installer}のように、\texttt{VSCodeUserSetup}で始まるインストーラをダブルクリックし、画面の指示に従ってインストールしてください。
ファイル名に含まれるバージョン番号は、画面写真と同じでなくても構いません。

\begin{figure}[H]
    \centering
    \includegraphics[width=.95\linewidth]{figures/chapter01/vscode-installer.png}
    \caption{ダウンロードしたVS Codeのインストーラを起動する}
    \label{fig:chapter01-vscode-installer}
\end{figure}

\begin{notebox}{インストール時の注意}
インストール中に「PATHに追加する」や「Codeで開く」という項目が表示された場合は、有効にしておくと便利です。授業中に設定が異なっていても大きな問題にはなりませんが、フォルダを右クリックしてVS Codeで開けるようになります。
\end{notebox}

インストール後は、Windowsの検索欄に\texttt{vscode}と入力し、「Visual Studio Code」を選ぶと起動できます（図\ref{fig:chapter01-start-vscode}）。
検索結果にインストーラも表示される場合がありますが、普段の作業で使うのはインストール済みの「Visual Studio Code」です。

\begin{figure}[H]
    \centering
    \includegraphics[width=.8\linewidth,height=.4\textheight,keepaspectratio]{figures/chapter01/start-vscode.png}
    \caption{Windowsの検索欄からVS Codeを起動する}
    \label{fig:chapter01-start-vscode}
\end{figure}

\subsection{Node.jsをインストールする}

Node.jsは、TypeScriptやp5.jsを使った教材プロジェクトを動かすために必要な実行環境です。Node.jsをインストールすると、\texttt{npm}というコマンドも一緒に使えるようになります。

\worktiming{PCごとに原則1回}

Webブラウザで\href{https://nodejs.org/}{Node.jsの公式サイト}を開き、図\ref{fig:chapter01-node-website}の「Node.jsを入手」を選びます。

\begin{figure}[H]
    \centering
    \includegraphics[width=.95\linewidth]{figures/chapter01/node-oficial.png}
    \caption{Node.jsの公式サイトからダウンロードページを開く}
    \label{fig:chapter01-node-website}
\end{figure}

ダウンロードページでは、Windows用で「LTS」と表示された版を選びます。LTS版は、長く安定して使うための版です。
図\ref{fig:chapter01-node-download}のように下へスクロールし、Windows Installer（\texttt{.msi}）をダウンロードします。
ページ上部にはDocker用のコマンドも表示されていますが、この教材では使いません。図の下部にあるインストーラのボタンを選んでください。

\begin{figure}[H]
    \centering
    \includegraphics[width=.9\linewidth,height=.52\textheight,keepaspectratio]{figures/chapter01/node-download.png}
    \caption{LTS版のWindowsインストーラ（\texttt{.msi}）を選ぶ}
    \label{fig:chapter01-node-download}
\end{figure}

教材で使うVite 8は、Node.js 20系では20.19.0以上、22系では22.12.0以上が必要です。
Node.js公式サイトの現在のLTS版インストーラはこの条件を満たします。
すでに古いNode.jsがインストールされている場合は、LTS版に更新してから使ってください。
インストーラでは標準設定のまま進め、インストールを最後まで完了してください。

インストールが終わったら、開いているVS CodeとPowerShellをいったん閉じてから、もう一度VS Codeを開きます。
コマンドを入力するには、VS Codeの「ターミナル」→「新しいターミナル」を選びます。
図\ref{fig:chapter01-open-terminal}のように「表示」→「ターミナル」から開くこともできます。英語表示の場合は「View」→「Terminal」です。

\begin{figure}[H]
    \centering
    \includegraphics[width=.95\linewidth]{figures/chapter01/start-termnial-on-vscode.png}
    \caption{VS Codeのメニューからターミナルを開く（英語表示の例）}
    \label{fig:chapter01-open-terminal}
\end{figure}

図\ref{fig:chapter01-powershell-terminal}では、画面下部にPowerShellのターミナルが開いています。
右上に\texttt{powershell}と表示されていることを確認します。別のものが開いた場合は、「+」の横のメニューから「PowerShell」を選んでください。
\texttt{PS C:\ldots>}の右側にコマンドを入力し、\texttt{Enter}で実行します。
この\texttt{PS C:\ldots>}や、サンプルの行番号は入力しません。ユーザー名やフォルダの場所は、自分のPCのものが表示されます。

\begin{figure}[H]
    \centering
    \includegraphics[width=.95\linewidth]{figures/chapter01/terninal-on-vsocde.png}
    \caption{画面下部に開いたPowerShellでコマンドを入力する}
    \label{fig:chapter01-powershell-terminal}
\end{figure}

新しく開いたターミナルで、次のコマンドを一行ずつ実行し、Node.jsとnpmが使えるか確認します。

\begin{listing}[htbp]
\begin{minted}{powershell}
node -v
npm -v
\end{minted}
\caption{Node.jsとnpmのバージョン確認}
\label{lst:chapter01-check-node}
\end{listing}

Node.jsでは\texttt{v}で始まる値、npmでは数字で始まるバージョン番号が表示されれば、Node.jsとnpmを使う準備ができています。なお、すでに古いNode.jsが入っている場合は、Vite 8の条件を満たすために更新が必要になることがあります。

\begin{notebox}{コマンドが認識されないとき}
\texttt{node}や\texttt{npm}が「コマンドとして認識されません」と表示されたときは、まずターミナルを閉じて開き直します。それでも直らないときはWindowsを再起動してください。再起動後も実行できない場合は、表示されたメッセージを添えて担当者に相談してください。
\end{notebox}

\subsection{Gitをインストールする}

Gitは、プログラムの変更を記録するための道具です。この教材では、GitHubに置かれている教材ファイルを取得するときにもGitを使います。

\worktiming{PCごとに原則1回}

Webブラウザで\href{https://git-scm.com/}{Gitの公式サイト}を開き、「Install for Windows」からWindows用のインストーラをダウンロードします。
図\ref{fig:chapter01-git-install}は、x64版のインストーラを選び、「ダウンロード」フォルダから起動するまでの流れです。
ARM版のWindowsを使っている場合は、ARM64版を選びます。どちらか分からない場合は、担当者に確認してください。
インストール項目はたくさんありますが、最初は基本的に標準設定のままで構いません。

インストールが終わったら、VS Codeのターミナルで次のコマンドを実行して、Gitが使えることを確認します。この章では、コマンドをPowerShellで実行する手順に統一します。

\begin{listing}[H]
\begin{minted}{powershell}
git --version
\end{minted}
\caption{Gitのバージョンを確認するコマンド}
\label{lst:chapter01-check-git}
\end{listing}

続いて、Gitに記録する名前とメールアドレスを設定します。次の\texttt{Taro Student}と\texttt{taro@example.com}は自分の情報に置き換えてください。メールアドレスには、授業の作業に使うことを許可されているものを指定します。

\begin{listing}[H]
\begin{minted}{powershell}
git config --global user.name "Taro Student"
git config --global user.email "taro@example.com"
\end{minted}
\caption{Gitに名前とメールアドレスを設定するコマンド}
\label{lst:chapter01-git-identity}
\end{listing}

\begin{figure}[H]
    \centering
    \includegraphics[width=.95\linewidth,height=.86\textheight,keepaspectratio]{figures/chapter01/git-install.png}
    \caption{Gitの公式サイトでWindows用のインストーラを選び、起動する}
    \label{fig:chapter01-git-install}
\end{figure}

\subsection{ファイルとフォルダをCLIで操作してみよう}
\label{sec:chapter01-cli-files}

教材を取得する前に、ターミナルでファイルやフォルダを操作する練習をします。
文字で命令を入力してコンピュータを操作する方法を、\textbf{CLI（Command Line Interface）}と呼びます。
入力する命令が\textbf{コマンド}で、その入力や結果の表示に使う画面がターミナルです。
この節では、VS Codeの統合ターミナルでPowerShellを使います。

\worktiming{最初に練習し、操作を忘れたときに見返す}

最も大切なのは、\textbf{現在どのフォルダにいるのかを意識すること}です。
コマンドは、特に場所を指定しなければ、現在のフォルダを基準にファイルを探したり作ったりします。
現在のフォルダは「カレントディレクトリ」とも呼びます。「ディレクトリ」と「フォルダ」は、ここではほぼ同じ意味です。

マウスでアイコンやメニューを選ぶ操作方法はGUIと呼びます。
VS CodeやWindowsのエクスプローラーでも、ファイルのコピーや名前変更はできます。
GUIを使ってはいけないわけではありません。Git、Node.js、npmを使うときにはターミナルを利用する機会が多いため、基本的なCLI操作にも慣れておきましょう。

\subsubsection*{練習する場所とターミナルを用意する}

まず、Windowsのエクスプローラーでドキュメントフォルダを開き、授業用の\texttt{work-textbook}フォルダを作ります。
すでに同名の授業用フォルダがある場合は、そのフォルダを使います。
VS Codeの「ファイル」→「フォルダーを開く」で、この\texttt{work-textbook}を開いてください。
信頼を確認する画面が出た場合は、自分で用意した授業用フォルダであることを確認してから信頼します。

続いて「ターミナル」→「新しいターミナル」を選びます。ターミナルの開き方と画面の見方は、図\ref{fig:chapter01-open-terminal}と図\ref{fig:chapter01-powershell-terminal}でも確認できます。
PowerShell以外が開いた場合は、ターミナル上部の「+」の横のメニューから「PowerShell」を選び、新しく開きます。
この節のコマンドは、そのPowerShellに一行ずつ入力し、\texttt{Enter}で実行してください。
画面の\texttt{PS C:\ldots>}のような入力待ちの表示や、掲載例の行番号は入力しません。

\begin{listing}[H]
\begin{minted}{powershell}
pwd
ls
\end{minted}
\caption{現在のフォルダと、その中のファイルを確認する}
\label{lst:chapter01-cli-location}
\end{listing}

\texttt{pwd}には現在のフォルダの場所が表示されます。末尾が\texttt{work-textbook}であることを確認します。
\texttt{ls}には、その中にあるファイルやフォルダの一覧が表示されます。
作ったばかりの空のフォルダなら、\texttt{ls}を実行しても一覧は表示されません。
違う場所が表示された場合は、\texttt{work-textbook}をVS Codeで開き直し、新しいターミナルで確認してください。

\begin{pointbox}{分からなくなったら現在地を確認する}
操作の途中で迷ったら、まず\texttt{pwd}で現在地を確認し、\texttt{ls}で周囲のファイルやフォルダを確認します。
VS Codeで開いているフォルダと、ターミナルで現在いるフォルダは、移動の操作によって異なる場合があります。
コマンドはすべて暗記する必要はありません。表を見ながら、この確認を習慣にしましょう。
\end{pointbox}

\begin{table}[htbp]
    \centering
    \caption{この節で使うPowerShellの基本操作}
    \label{tab:chapter01-cli-commands}
    \begin{tabular}{p{0.30\linewidth}p{0.60\linewidth}}
        \toprule
        コマンドの例 & 意味 \\
        \midrule
        \texttt{pwd} & 現在いるフォルダを確認する \\
        \texttt{ls} & 現在のフォルダの内容を表示する \\
        \texttt{cd practice} & 現在のフォルダ内にある\texttt{practice}へ移動する \\
        \texttt{cd ..} & 1つ上のフォルダへ移動する \\
        \texttt{mkdir practice} & \texttt{practice}フォルダを新しく作る \\
        \texttt{cp memo.txt test.txt} & 元を残して\texttt{test.txt}へコピーする \\
        \texttt{mv test.txt my-test.txt} & \texttt{test.txt}の名前を\texttt{my-test.txt}へ変更する \\
        \texttt{rm my-test.txt} & \texttt{my-test.txt}を削除する \\
        \bottomrule
    \end{tabular}
\end{table}

表の\texttt{practice}や\texttt{memo.txt}は、操作するフォルダ名やファイル名です。
名前はコマンドの後に半角の空白を入れて指定します。名前や場所に空白が含まれる場合は、\texttt{cd "My Folder"}のように引用符で囲みます。
以下では入力しやすい英数字の名前で練習します。

\subsubsection*{フォルダを作り、その中へ移動する}

\texttt{work-textbook}の中に、練習専用の\texttt{practice}を作って移動します。
すでに練習用の\texttt{practice}がある場合は、\texttt{mkdir practice}は繰り返さず、\texttt{cd practice}から始めてください。

\begin{listing}[H]
\begin{minted}{powershell}
mkdir practice
cd practice
pwd
ls
\end{minted}
\caption{練習用フォルダを作り、その中へ移動する}
\label{lst:chapter01-cli-practice}
\end{listing}

\texttt{mkdir}はフォルダを作るだけで、その中へ移動するわけではありません。
\texttt{cd practice}で移動してから\texttt{pwd}を実行すると、場所の末尾が\texttt{work-textbook\textbackslash practice}になります。
\texttt{ls}で練習用フォルダの内容を確認しましょう。

\subsubsection*{ファイルをコピーし、名前を変更する}

教材はまだ取得していないので、ここでは自分で短い練習用ファイルを用意します。
VS Codeのエクスプローラーで\texttt{practice}を選び、その中に\texttt{memo.txt}を作ります。
「CLIの練習」のような一行を書き、\texttt{Ctrl + S}で保存してください。
ターミナルに戻り、\texttt{pwd}と\texttt{ls}で\texttt{practice}の中に\texttt{memo.txt}があることを確認します。

\begin{listing}[H]
\begin{minted}{powershell}
cp memo.txt test.txt
ls
mv test.txt my-test.txt
ls
\end{minted}
\caption{練習用ファイルをコピーし、コピーの名前を変更する}
\label{lst:chapter01-cli-copy-rename}
\end{listing}

\texttt{cp}では、先にコピー元、次にコピー先を指定します。
最初の\texttt{ls}で、元の\texttt{memo.txt}とコピーの\texttt{test.txt}の両方があることを確認してください。
次の\texttt{mv}では、現在の名前と変更後の名前を指定しています。
最後の\texttt{ls}では\texttt{test.txt}がなくなり、\texttt{memo.txt}と\texttt{my-test.txt}が表示されます。
繰り返し練習する場合は、同名のコピー先が残っていないかを先に確認し、残っている場合は別の名前を使ってください。

\texttt{mv}には、名前変更のほか、ファイルを別のフォルダへ移動する使い方もあります。
移動先に既存のフォルダを指定すると、その中へファイルが移ります。
\texttt{cp}と違い、移動後は元の場所に同じファイルは残りません。

\subsubsection*{コピーだけを削除し、元のフォルダへ戻る}

最後に、練習で作ったコピーの\texttt{my-test.txt}だけを削除します。
元の\texttt{memo.txt}や、ほかのファイルは削除しません。

\begin{notebox}{\texttt{rm}はファイルを削除するコマンド}
\texttt{rm}で削除したファイルは、Windowsのエクスプローラーでの操作と違って、ごみ箱には入らず削除されます。
ごみ箱から戻せるとは考えないでください。
削除前には必ず\texttt{pwd}で現在地を、\texttt{ls}で対象のファイルを確認します。
次の例では\texttt{practice}の中にいて、削除する名前が練習用の\texttt{my-test.txt}であることを確かめてから実行してください。
\end{notebox}

\begin{listing}[H]
\begin{minted}{powershell}
pwd
ls
rm my-test.txt
ls
cd ..
pwd
ls
\end{minted}
\caption{練習用のコピーを削除し、1つ上のフォルダへ戻る}
\label{lst:chapter01-cli-delete-return}
\end{listing}

削除後の\texttt{ls}で、\texttt{memo.txt}は残り、\texttt{my-test.txt}だけがなくなったことを確認します。
\texttt{cd ..}の\texttt{..}は「1つ上のフォルダ」を表します。
戻った後の\texttt{pwd}では\texttt{work-textbook}、\texttt{ls}では\texttt{practice}が確認できます。
\texttt{practice}と\texttt{memo.txt}は、次回の練習用に残しておきましょう。

この後は、\texttt{work-textbook}で教材を取得し、取得した教材の中へ移動します。
\texttt{npm install}や\texttt{npm run dev}では、さらに\texttt{package.json}がある\texttt{typescript-p5}へ移動して実行します。
コマンドの前に現在地を確認することが、これらの作業でも役に立ちます。

\subsection{教材をGitHubから取得する}

教材のファイルはGitHubから取得します。GitHubは、プログラムや教材ファイルを保存・共有するためのサービスです。
この教材のURLは\url{https://github.com/m-riho/typescript-p5js-programming-textbook}です。

\worktiming{教材をこのPCへ取得するときに原則1回}

前節で用意した\texttt{work-textbook}に教材を取得します。
\texttt{pwd}で現在地が\texttt{work-textbook}であることを、\texttt{ls}でその内容を確認してから\texttt{git clone}を実行してください。
別の場所にいる場合は、VS Codeで\texttt{work-textbook}を開き、新しいPowerShellターミナルで確認します。

\begin{listing}[htbp]
\begin{minted}{powershell}
pwd
ls
git clone https://github.com/m-riho/typescript-p5js-programming-textbook.git
ls
\end{minted}
\caption{現在地を確認し、教材のリポジトリを取得するコマンド}
\label{lst:chapter01-git-clone}
\end{listing}

上のURLは、この教材の学生向け公開リポジトリの正式なURLです。\texttt{git clone}は通常、最初に教材を取得するときに一度だけ実行します。
実行が終わると、最後の\texttt{ls}に、練習用の\texttt{practice}と取得した\texttt{typescript-p5js-programming-textbook}が表示されます。
\texttt{git clone}は教材フォルダを作りますが、その中へは移動しません。次の節で、取得した教材フォルダをVS Codeで開きます。
すでに教材フォルダがある場合は、\texttt{git clone}を繰り返さず、次の節へ進んでください。

図\ref{fig:chapter01-get-material}は、コマンドの入力と、取得した教材フォルダをエクスプローラーで確認する様子です。
画像中の\texttt{git clone}の入力例は\texttt{Documents}で実行していますが、\textbf{本書の手順では、先に\texttt{work-textbook}へ移動してから実行します}。
\texttt{mkdir work-textbook}だけでは、その中へ移動しないことにも注意してください。
操作には画像中のコマンドではなく、サンプル\ref{lst:chapter01-git-clone}を使います。

\begin{figure}[H]
    \centering
    \includegraphics[width=\linewidth,height=.72\textheight,keepaspectratio]{figures/chapter01/get-material-from-git.png}
    \caption{教材取得の操作画面の例（実行場所は本文の手順に合わせる）}
    \label{fig:chapter01-get-material}
\end{figure}

\begin{notebox}{フォルダ名に注意する}
Windowsでは、フォルダ名に日本語や空白が含まれていても多くの場合は動作します。ただし、コマンド入力に慣れるまでは、\texttt{work-textbook}のような英数字だけのフォルダ名にしておくとトラブルを減らせます。
\end{notebox}

\subsection{VS Codeで教材フォルダを開く}

教材を取得したら、VS Codeでそのリポジトリのルートフォルダを開きます。VS Codeのメニューから「ファイル」→「フォルダーを開く」を選び、\texttt{git clone}で作られた教材フォルダを指定します。

\worktiming{作業を始めるたび（初回準備でも実行）}

フォルダを開いたら、VS Codeの画面左側にあるエクスプローラーで\texttt{README.md}と\texttt{typescript-p5/}が表示されることを確認します。\texttt{typescript-p5/}を展開すると、\texttt{listings/}、\texttt{workspace/}、\texttt{src/}などが見つかります。プログラムを書くときは、この教材フォルダを開いた状態で作業します。
学生向けリポジトリには、サンプルプログラムと実行に必要なファイルを用意しています。教材PDFはGitHubのReleaseから取得できます。本文を編集・組版するためのLaTeXソースは別の教員向けリポジトリに置いてあり、授業の準備では取得する必要はありません。

初めて開くフォルダでは、ワークスペースを信頼するか尋ねる画面が表示されることがあります。授業で配布されたリポジトリを開いていることを確認してから、信頼する選択肢を選びます。統合ターミナルは、VS Codeのメニューから「ターミナル」→「新しいターミナル」を選ぶと開けます。

\begin{pointbox}{自分のプログラムを保存する場所}
授業中の練習や自由制作で作るTypeScriptファイルは、\texttt{typescript-p5/}の中の\texttt{workspace/}以下に保存します。
この章で\texttt{workspace/}と書く場合は、この作業用フォルダを指します。\texttt{listings/}には、変更せずに残しておく教材のサンプルがあります。
演習問題・課題のプログラムは、\texttt{workspace/exercises/}に保存してください。
ファイル名は\texttt{章番号-問題番号.ts}とし、番号をそれぞれ2桁にそろえます。
たとえば、第2章の問題1は\texttt{02-01.ts}、第3章の問題5は\texttt{03-05.ts}、第10章の問題3は\texttt{10-03.ts}です。
ファイル名で並べると章・問題の順番になり、課題を見つけやすくなります。
\end{pointbox}

教材本体の\texttt{typescript-p5/src/}や設定ファイルは、教材や担当教員の指示がある場合にだけ変更します。
この章の後半では、\texttt{typescript-p5/index.html}の実行先を変更する操作を説明します。
\texttt{workspace/}へ保存しただけでは、実行するファイルとして選択されません。\texttt{index.html}でそのファイルを指定すると、保存したプログラムを直接実行できます。\texttt{src/main.ts}へ内容をコピーする必要はありません。

\subsection{必要なライブラリをインストールする}

教材プロジェクトには、p5.jsやTypeScriptなど、プログラムを動かすために必要なライブラリが含まれています。ただし、GitHubから取得した直後は、ライブラリ本体はまだ手元に入っていません。

\worktiming{教材プロジェクトごとに原則1回。必要なライブラリが変わった場合は再実行}

VS Codeで教材リポジトリのルートフォルダを開いた状態で、メニューから「ターミナル」→「新しいターミナル」を選びます。画面下部にターミナルが開いたら、次のコマンドを上から順に実行します。

\begin{listing}[htbp]
\begin{minted}{powershell}
cd typescript-p5
pwd
ls package.json
npm install
\end{minted}
\caption{実行用プロジェクトで必要なライブラリをインストールするコマンド}
\label{lst:chapter01-npm-install}
\end{listing}

\texttt{package.json}は\texttt{typescript-p5/}の中にあります。\texttt{pwd}で現在の場所を確かめ、\texttt{ls package.json}でファイルが表示されることを確認してから、必ず\texttt{typescript-p5/}で\texttt{npm install}を実行します。
\texttt{ls}の後にファイル名を指定すると、そのファイルに絞って確認できます。
\texttt{npm install}は、教材プロジェクトに必要なライブラリをまとめてインストールするコマンドです。完了まで少し時間がかかることがあります。

実行すると\texttt{node\_modules/}フォルダが自動的に作られます。また、\texttt{package-lock.json}には、インストールしたライブラリのバージョンが記録されます。表示に警告行が含まれていても、その内容を確認してください。一方、出力の終わり近くに\texttt{npm error}を含む行がないか確認し、その前後に表示されたメッセージを読んで、インストールが成功していない場合の原因を確認します。

\begin{pointbox}{\texttt{npm install}は最初の1回}
\texttt{npm install}は、基本的には教材プロジェクトを最初に準備するときに1回実行します。
別のPCで作業するときや、教材プロジェクトに必要なライブラリが変わったときには、もう一度実行してください。
\end{pointbox}

\begin{notebox}{エラーが出たとき}
\texttt{npm install}でエラーが出た場合は、まずターミナルが\texttt{typescript-p5/}を開いているか確認してください。ターミナルに表示されている場所が違う場合、\texttt{package.json}が見つからず失敗します。
\end{notebox}

\section{プログラムを動かす編}
ここでは、教材の準備が終わったあと、授業や制作で繰り返し使う基本操作を確認します。作業を始めるたびに、この編の早見表を使ってください。

\begin{chapterphasebox}{授業や制作のたびに使う}{orange}
VS Codeで教材を開き、開発サーバを起動し、プログラムを編集します。Gitのコミットは毎回機械的に行うのではなく、意味のある作業の節目で行います。
\end{chapterphasebox}

\begin{pointbox}{毎回の作業早見表}
\begin{enumerate}
    \item VS Codeで教材フォルダを開く
    \item ターミナルで\texttt{typescript-p5/}へ移動し、\texttt{pwd}と\texttt{ls package.json}で確認する
    \item \texttt{npm run dev}で開発サーバを起動する
    \item ブラウザで表示を確認する
    \item \texttt{index.html}で実行するファイルを選ぶ
    \item 自分のプログラムを\texttt{workspace/}以下に作り、直接編集して保存する
    \item 意味のある作業の節目でGitへコミットする
    \item 作業を終えるときに\texttt{Ctrl + C}で開発サーバを止める
\end{enumerate}
\end{pointbox}

\subsection{開発サーバを起動する}

ライブラリのインストールが終わったら、開発サーバを起動します。開発サーバは、作成中のプログラムをブラウザで確認するための仕組みです。

\worktiming{作業を始めるたび}

VS Codeのターミナルで、\texttt{pwd}と\texttt{ls package.json}を使って現在地を確認します。
\texttt{typescript-p5/}にいて、\texttt{package.json}が表示されたら、次のコマンドを実行します。

\begin{listing}[htbp]
\begin{minted}{powershell}
npm run dev
\end{minted}
\caption{開発サーバを起動するコマンド}
\label{lst:chapter01-npm-run-dev}
\end{listing}

実行すると、ターミナルに\texttt{http://localhost:5173/}や\texttt{http://localhost:5174/}のようなLocal URLが表示されます。URLは\texttt{Ctrl}を押しながらクリックするか、コピーしてブラウザで開きます。5173番のポートがすでに使われている場合は、5173以外の番号になることがありますが、正常な動作です。

\texttt{npm run dev}を実行すると、教材プロジェクトではViteという開発用のWebサーバが起動します。
ブラウザで表示されたURLを開いたときに、最初に読み込まれるのは\texttt{typescript-p5/index.html}です。
\texttt{index.html}は、Webページの入口になるHTMLファイルです。

\begin{listing}[htbp]
\begin{minted}{html}
<body>
  <div id="app"></div>
    <script type="module"
            src="/listings/chapter01/first-sketch.ts"></script>
</body>
\end{minted}
\caption{\texttt{index.html}から最初のサンプルを読み込む部分}
\label{lst:chapter01-index-html-main-ts}
\end{listing}

このHTMLでは、\texttt{div}要素でページ内に場所を用意し、\texttt{script}要素でTypeScriptの入口になるファイルを読み込んでいます。
読み込む先として指定されているのが、\texttt{/listings/chapter01/first-sketch.ts}です。
先頭の\texttt{/}は、Viteが公開するフォルダ\texttt{typescript-p5/}を表します。実際のファイルは、この中の\texttt{listings/chapter01/first-sketch.ts}です。\texttt{src}属性に\texttt{/typescript-p5/}と書き足す必要はありません。
ブラウザでページを開くと、まず\texttt{index.html}が読み込まれ、そこに指定したTypeScriptファイルが読み込まれます。そのファイルの中でp5.jsを取り込み、スケッチを開始すると、画面に図形が表示されます。

\begin{pointbox}{localhostとは}
\texttt{localhost}は、自分のPC自身を表す名前です。インターネット上に公開しているわけではなく、自分のPCの中で動いている開発用のページを見ています。
\end{pointbox}

\begin{notebox}{開発サーバを止める}
\texttt{npm run dev}で開発サーバを起動している間、そのターミナルは開発サーバ用に使われます。
別のコマンドを実行したい場合は、VS Codeで新しいターミナルを開いてください。
開発サーバを止めるときは、ターミナルで\texttt{Ctrl + C}を押します。PowerShellが確認を求めた場合だけ、\texttt{Y}を入力します。
\end{notebox}

\subsection{プログラムを実行してみよう}

ブラウザで開いたページに図形が表示されれば、TypeScript + p5.jsのプログラムが動いています。配布時に選択されているファイルは、\texttt{typescript-p5/}の中の\texttt{listings/chapter01/first-sketch.ts}です。ここでは、最初のスケッチの形だけを見ておきます。

\worktiming{開発サーバを起動したあとに確認}

\begin{listing}[htbp]
\inputminted{ts}{typescript-p5/listings/chapter01/first-sketch.ts}
\caption{最初のp5.jsスケッチ（\texttt{chapter01/first-sketch.ts}）}
\label{lst:chapter01-first-sketch}
\end{listing}

\begin{pointbox}{最初は定型として使う}
このスケッチには、これから学ぶ構文がいくつか含まれています。最初からすべてを理解したり、暗記したりする必要はありません。当面は、\texttt{import}から\texttt{new p5(sketch);}までをスケッチを動かすための定型として使います。まずは\texttt{p.background(...)}や\texttt{p.circle(...)}の数値を書き換え、表示がどのように変わるかを確かめましょう。
\end{pointbox}

最初のスケッチには、JavaScriptの書き方、TypeScriptで追加された書き方、p5.jsが提供する機能が一緒に現れます。これらは、次の表のように分けて考えることができます。それぞれの詳しい意味は、必要になる章で順に説明します。

\begin{table}[htbp]
    \centering
    \caption{最初のスケッチに含まれる3種類の書き方}
    \label{tab:chapter01-javascript-typescript-p5}
    \begin{tabular}{p{0.24\linewidth}p{0.27\linewidth}p{0.39\linewidth}}
        \toprule
        種類 & 例 & この教材での意味 \\
        \midrule
        JavaScriptの書き方
            & \texttt{const}、関数、\texttt{=}
            & 値に名前を付け、処理を関数にまとめ、値や処理を代入する書き方です。 \\
        TypeScriptの書き方
            & \texttt{: p5}、\texttt{: void}
            & 引数や戻り値がどのような型なのかを示す型注釈です。 \\
        p5.jsの機能
            & \texttt{p.setup}、\texttt{p.circle}
            & 実行したい処理を登録したり、画面に図形を描いたりするための機能です。 \\
        \bottomrule
    \end{tabular}
\end{table}

このプログラムの2行目の\texttt{import p5 from "p5";}では、p5.jsの機能を取り込んでいます。
5行目から15行目では、\texttt{sketch}という関数の中で、p5.jsの初期化処理(\texttt{setup})と描画処理(\texttt{draw})を定義しています。
つまり、p5.jsに、どのようなプログラムを動かしてほしいかを登録していることになります。
最後に\texttt{new p5(sketch)}でp5.jsのインスタンスを作成し、\texttt{sketch}関数を渡すことで、プログラムが実行されます。
つまり、この設計図を使って、実際にp5.jsを動かし始めることになります。

このプログラムでは、\texttt{p.createCanvas(400, 300)}で幅400、高さ300の描画領域を作り、\texttt{p.circle(200, 150, 80)}で円を描いています。
p5.jsの機能は、スケッチを表す変数を通して呼び出します。本書ではその変数を主に\texttt{p}とし、\texttt{p.createCanvas(...)}や\texttt{p.circle(...)}のように書きます。

\begin{figure}[htbp]
    \centering
    \includegraphics[width=0.5\textwidth]{figures/chapter01/browser-first-sketch.png}
    \caption{最初のスケッチの表示例}
    \label{fig:chapter01-first-sketch}
\end{figure}

5行目の\texttt{sketch}関数の引数名は、必ずしも\texttt{p}である必要はありません。
7行目や11行目で使っている\texttt{p.setup}や\texttt{p.draw}の\texttt{p}も、引数名を変えた場合は同じ名前にそろえます。
たとえば、引数名を\texttt{myP5}にした場合は、\texttt{myP5.createCanvas(...)}や\texttt{myP5.circle(...)}のように書きます（Listing~\ref{lst:chapter01-first-sketch-myP5}）。
慣習的に\texttt{p}という名前を使うのが慣習のようですので、この資料では\texttt{p}を使うことにします。

\begin{listing}[htbp]
\inputminted{ts}{typescript-p5/listings/chapter01/first-sketch-myP5.ts}
\caption{最初のp5.jsスケッチの\texttt{p}を\texttt{myP5}に変更した例（\texttt{chapter01/first-sketch-myP5.ts}）}
\label{lst:chapter01-first-sketch-myP5}
\end{listing}

Listing~\ref{lst:chapter01-first-sketch-myP5}を試すには、\texttt{typescript-p5/index.html}を開き、既存の\texttt{script}要素の\texttt{src}属性を次のように変更して保存します。\texttt{script}要素を追加するのではなく、実行先の1か所を書き換えます。

\begin{listing}[htbp]
\begin{minted}{html}
<script type="module"
        src="/listings/chapter01/first-sketch-myP5.ts"></script>
\end{minted}
\caption{\texttt{myP5}を使ったサンプルに実行先を切り替える}
\label{lst:chapter01-select-myP5}
\end{listing}

\subsection{プログラムを書き換えてみよう}

プログラムが動いたら、自分の作業用ファイルを作って少しだけ書き換えてみます。最初は、円の大きさ、円の位置、背景の色を変えるだけで十分です。元のサンプルは残し、自分のファイルを編集します。

\worktiming{作業中}
VS Codeのエクスプローラーで\texttt{typescript-p5/}の中の\texttt{workspace/}を選び、\texttt{first-sketch.ts}という新しいファイルを作ります。\texttt{listings/chapter01/first-sketch.ts}の内容全体を新しいファイルへコピーして保存してください。ここでは\texttt{myP5}の例ではなく、\texttt{p}を使う最初のサンプルをもとにします。

続いて\texttt{index.html}の実行先を、Listing~\ref{lst:chapter01-select-workspace}のように変更して保存します。それ以降は\texttt{workspace/first-sketch.ts}だけを編集します。変更したら\texttt{Ctrl + S}で保存し、ブラウザで結果を確認しましょう。

\begin{listing}[htbp]
\begin{minted}{html}
<script type="module" src="/workspace/first-sketch.ts"></script>
\end{minted}
\caption{自分の作業用ファイルを直接実行する}
\label{lst:chapter01-select-workspace}
\end{listing}

たとえば、\texttt{p.circle(200, 150, 80)}の最後の数値を\texttt{120}に変えると、円が大きくなります。また、\texttt{p.background(240)}の数値を変えると、背景の明るさが変わります。

\begin{listing}[htbp]
\begin{minted}{ts}
p.background(220);
p.circle(200, 150, 120);
\end{minted}
\caption{背景と円の大きさを変更する例}
\label{lst:chapter01-edit-circle}
\end{listing}

課題を作る場合も、\texttt{workspace/exercises/}に指定の名前のファイルを作ります。たとえば、この章の演習問題7なら\texttt{01-07.ts}とし、\texttt{index.html}の実行先を\texttt{/workspace/exercises/01-07.ts}にします。次回も同じファイルを指定すれば、保存したプログラムをそのまま実行できます。\texttt{src/main.ts}へコピーしたり、2つのファイルの内容をそろえたりする作業は不要です。

別の章のサンプルを試す場合は、キャプションにあるディレクトリ名とファイル名を確認し、\texttt{index.html}の実行先を変更します。第2章の点と線の例は、Listing~\ref{lst:chapter01-copy-sample}のように指定します。

\begin{listing}[htbp]
\begin{minted}{html}
<script type="module"
        src="/listings/chapter02/points-and-lines.ts"></script>
\end{minted}
\caption{第2章のサンプルを直接実行する指定}
\label{lst:chapter01-copy-sample}
\end{listing}

指定したファイルそのものが実行されます。保存した後は、通常はViteを再起動しなくてもブラウザの表示が更新されます。画像、JSON、音を使うサンプルでは、それぞれの章で説明する素材ファイルも必要です。

\begin{notebox}{単独実行用のサンプルと説明用のコード片}
後半の章には、関数やクラスの一部分だけを示すファイルもあります。すべての\texttt{.ts}ファイルが単独で動くわけではありません。\texttt{typescript-p5/listings/README.md}に、単独では実行しないファイルと、コンソールにだけ結果を表示するファイルをまとめています。まずは、この章で指定したサンプルで操作を確かめてください。
\end{notebox}

書き換えたあとに期待どおりに動かないときは、まずVS Codeの赤い波線を確認します。次にターミナルの表示を確認し、最後にブラウザの表示を確認します。

型の間違いを調べるときは、\texttt{typescript-p5/}にいる別のターミナルで\texttt{npm run check}を実行します。\texttt{workspace/}内のファイルも検査されます。別の未完成の課題にエラーがある場合も表示されるため、メッセージ中のファイル名を確認しましょう。\texttt{npm run dev}は表示を試すための処理で、型チェックとは別です。

\subsection{Gitで保存してみよう}

Gitでは、ファイルの変更を記録として保存できます。この記録をコミットと呼びます。コミットを作っておくと、あとで「どの時点で何を変更したか」を確認できます。この節では、開発サーバ用とは別にターミナルを開き、教材リポジトリのルートフォルダでGitコマンドを実行します。

\worktiming{意味のある作業の節目}

まず、現在の変更状態を確認します。

\begin{listing}[htbp]
\begin{minted}{powershell}
git status
\end{minted}
\caption{Gitで変更状態を確認するコマンド}
\label{lst:chapter01-git-status}
\end{listing}

変更したファイルが表示されたら、次のように保存対象へ追加します。先ほど作成した\texttt{workspace/first-sketch.ts}と、実行先を記載した\texttt{index.html}を一緒に記録します。

\begin{listing}[htbp]
\begin{minted}{powershell}
git add typescript-p5/index.html
git add typescript-p5/workspace/first-sketch.ts
\end{minted}
\caption{変更したファイルを保存対象へ追加するコマンド}
\label{lst:chapter01-git-add}
\end{listing}

この節ではGitコマンドを教材リポジトリのルートフォルダで実行します。一方、\texttt{npm}のコマンドは\texttt{typescript-p5/}で実行します。コマンドごとに実行する場所を確認しましょう。

課題を記録するときは、2行目のパスを実際に保存した課題ファイルのパスに置き換えます。Gitはリポジトリのルートで操作するので、\texttt{typescript-p5/}から書きます。
たとえば、この章の演習問題7では\texttt{typescript-p5/workspace/exercises/01-07.ts}です。\texttt{git status}で対象のファイル名を確認してから追加してください。

最後に、変更内容を説明する短いメッセージを付けてコミットします。

\begin{listing}[htbp]
\begin{minted}{powershell}
git commit -m "最初のスケッチを変更"
\end{minted}
\caption{Gitで変更を保存するコマンド}
\label{lst:chapter01-git-commit}
\end{listing}

\begin{notebox}{コミットメッセージ}
コミットメッセージは、あとで見返したときに変更内容が分かるように書きます。最初は短い日本語で構いません。
\end{notebox}

\section{よくある間違い}

環境構築では、プログラムの文法とは別のところでつまずくことがあります。エラーが出たときは、あわてずに表示されたメッセージと作業しているフォルダを確認しましょう。

\begin{mistakebox}{初回のコマンドを授業のたびに繰り返す}
\texttt{git clone}と\texttt{npm install}は、通常は教材を準備するときに1回実行します。すでに教材フォルダと\texttt{node\_modules/}がある場合は、毎回繰り返す必要はありません。一方、\texttt{npm run dev}は、プログラムを動かす作業を始めるたびに実行します。
\end{mistakebox}

\begin{mistakebox}{コマンドを実行する場所が違う}
\texttt{npm install}や\texttt{npm run dev}は、\texttt{package.json}がある\texttt{typescript-p5/}で実行します。
VS Codeで教材全体を開いていても、ターミナルの現在地が同じとは限りません。\texttt{pwd}と\texttt{ls package.json}で、コマンドを実行する場所を確認してください。
\end{mistakebox}

\begin{mistakebox}{開発サーバを止めてしまう}
\texttt{npm run dev}を実行しているターミナルを閉じると、開発サーバも止まります。
作業中はそのターミナルを開いたままにしておきます。
止めたいときは、ターミナルで\texttt{Ctrl + C}を押します。
\end{mistakebox}

\begin{mistakebox}{Gitの保存とファイル保存を混同する}
VS Codeでファイルを保存することと、Gitでコミットすることは別の操作です。ファイル保存は現在の内容をディスクに書き込む操作で、コミットは変更の記録を作る操作です。
\end{mistakebox}

\section{まとめ}

この章は、はじめに編、準備編、プログラムを動かす編の3つに分かれています。はじめに編は最初に読んで必要なときに見返し、準備編はPCや教材ごとに原則1回行います。

CLI操作では、\texttt{pwd}で現在地を確認し、\texttt{ls}で内容を見て、\texttt{cd}で移動することを基本にします。
コマンドをすべて暗記するよりも、どのフォルダの何を操作するのかを確かめる習慣を身につけましょう。

プログラムを動かす編は、授業や制作のたびに使います。\texttt{npm run dev}で開発サーバを起動し、\texttt{index.html}で実行先を選び、自分の\texttt{workspace/}内のファイルを編集します。\texttt{git clone}や\texttt{npm install}は毎回繰り返さず、Gitのコミットは意味のある作業の節目で行いましょう。次章からは、p5.jsを使って図形を描く方法を少しずつ学んでいきます。

\section{演習問題}

この章の演習では、環境構築と最初の実行を確認します。すべてを暗記する必要はありません。コマンドを見ながら、同じ操作を自分のPCで再現できることを目指しましょう。

プログラムを作成・変更する問題では、内容を\texttt{workspace/exercises/}へ保存します。問題7は\texttt{01-07.ts}、問題10は\texttt{01-10.ts}とします。コマンドの確認だけの問題について、TypeScriptファイルを作る必要はありません。

\begin{enumerate}
    \item \texttt{git clone}、\texttt{npm install}、\texttt{npm run dev}、\texttt{git commit}を、「原則1回」「作業開始時」「作業の節目」に分けてください。
    \item VS Codeを起動し、教材フォルダを開いてください。
    \item VS CodeでPowerShellのターミナルを開き、\texttt{pwd}と\texttt{ls}で現在地と内容を確認してください。
    続いて\ref{sec:chapter01-cli-files}節を見ながら、授業用の\texttt{work-textbook}で\texttt{practice}の作成・移動、\texttt{memo.txt}のコピー・名前変更・コピーの削除を練習してください。
    すでにある練習用フォルダや元ファイルはそのまま使い、削除前には対象を確認します。
    練習後は、教材リポジトリのフォルダをVS Codeで開き直し、新しいターミナルで現在地を確認してください。
    \item \texttt{node -v}と\texttt{npm -v}を実行し、バージョン番号が表示されることを確認してください。
    \item ターミナルで\texttt{typescript-p5/}へ移動し、\texttt{pwd}と\texttt{ls package.json}で現在地を確認してください。
    その中に\texttt{node\_modules/}がなければ、\texttt{npm install}を実行してください。すでにある場合は、再実行する必要はありません。
    \item \texttt{npm run dev}を実行し、ブラウザで教材ページを開いてください。
    \item 最初のスケッチで、円の大きさと背景の明るさを変更してください。
    \item \texttt{git status}で変更を確認し、\texttt{git add}と\texttt{git commit}で作業の節目を記録してください。
    \item \texttt{index.html}の実行先を\texttt{/listings/chapter01/first-ballon-game.ts}にして、表示された風船をクリックしてみてください。このプログラムでは、本書の後半までに学ぶ画像の読み込み、配列、繰り返し、関数、クラスなどを使っています。現時点でプログラムの内容を理解する必要はありません。風船が動き、クリックで消せることを確認できれば十分です。
    \item 挑戦問題として、最初のスケッチの円の位置、大きさ、背景色をすべて変更し、自分なりの最初の画面を作ってください。
\end{enumerate}
